% !TEX root = ../notes_template.tex
\chapter{Lumbopelvic}\label{chp:lumbopelvic}

\section*{Objectives}
\begin{itemize}
  \item Apply regional interdependence and SFMA logic to lumbopelvic movement assessment.
  \item Perform manual muscle testing for trunk flexion and extension using graded criteria.
  \item Evaluate core motor control using the leg-lowering test with angle documentation.
  \item Identify primary muscles and innervation involved in trunk stability and motion.
  \item Acquire and interpret RUSI images of the transversus abdominis and lumbar multifidi.
  \item Differentiate activation patterns, thickness changes, and asymmetries relevant to clinical core dysfunction.
\end{itemize}

\section*{Functional Integration – Lumbopelvic Control}

The lumbopelvic region serves as a central force transducer and stabilizer between the lower extremities and thorax. It plays a crucial role in both static postural alignment and dynamic movement coordination. Closed-chain motions of the pelvis represent movements of the lumbar spine, and dysfunction here often manifests distally (e.g., hip or knee) or proximally (e.g., scapular instability, rib restriction). Functional control in this region is critical for injury prevention and efficient movement.

\subsection*{SFMA Top Tier Patterns with Lumbopelvic Focus}

\begin{itemize}
  \item \textbf{Multisegmental Flexion (MSF)} — Observe lumbar contribution to forward bend; compensatory hip strategies may indicate stiffness or motor control deficit.
  \item \textbf{Multisegmental Extension (MSE)} — Assess lumbar extension symmetry; anterior pelvic tilt or gluteal inhibition may alter pattern.
  \item \textbf{Multisegmental Rotation (MSR)} — View spinal dissociation and segmental contribution; asymmetries suggest regional mobility or timing deficits.
  \item \textbf{Single Leg Stance (SLS)} — Monitor pelvic stability and trunk sway; frontal plane drift may reflect core or hip deficits.
  \item \textbf{Arms Down Squat (ADS)} — Evaluate lumbar control under load; segmental flexion lag may reveal dysfunction.
\end{itemize}

\section*{Trunk Movements: Muscles and Innervation}

\begin{table}[H]
\centering
\captionsetup{justification=raggedright, singlelinecheck=false}
\begin{tabular}{|p{4cm}|p{5cm}|p{4cm}|p{2.5cm}|}
\hline
\textbf{Movement} & \textbf{Muscle} & \textbf{Innervation} & \textbf{Nerve Root(s)} \\
\hline
\multirow{2}{*}{Trunk Flexion}
  & Rectus Abdominis & Intercostal & T6–T12 \\
  & External Oblique & Intercostal, Subcostal & T7–T12 \\
\hline
\multirow{3}{*}{Trunk Extension}
  & Erector Spinae Group & Dorsal Rami & Multisegmental \\
  & Multifidus & Dorsal Rami & Multisegmental \\
  & Quadratus Lumborum & Ventral Rami & T12–L3 \\
\hline
\multirow{3}{*}{Lateral Flexion \& Rotation}
  & Internal Oblique & Intercostal \& Iliohypogastric & T7–L1 \\
  & External Oblique & Intercostal \& Subcostal & T7–T12 \\
  & Quadratus Lumborum & Ventral Rami  & T12–L3 \\
\hline
\multirow{2}{*}{Trunk Stabilization}
  & Transversus Abdominis & Intercostal, Iliohypogastric, Ilioinguinal & T7–L1 \\
  & Multifidus & Dorsal Rami & Multisegmental \\
\hline
\end{tabular}
\caption{Primary Trunk Movements with Neuromuscular Innervation}
\end{table}

\subsection*{MMT Activities – Lumbopelvic Region}

Unlike traditional manual muscle tests that isolate limb actions against gravity or resistance, lumbopelvic MMTs assess global trunk patterns—often in a closed-chain or bodyweight-supported context. These tests reflect the functional nature of the core: to stabilize, transfer load, and coordinate movement across multiple segments rather than generate isolated joint motion.

Grading for trunk flexion and extension is based on standardized positioning and patient control relative to gravity, not therapist-applied resistance. This approach emphasizes movement quality, segmental control, and compensatory strategies, aligning with functional anatomy and the demands of daily movement.

The following tests are an approach for assessing anterior and posterior core function, as well as integrated trunk control through leg-lowering. Interpret each test not just for strength but for the presence of motor control deficits, endurance limitations, or regional asymmetries.

\begin{labactivitybox}
\textbf{Activity: MMT – Trunk Flexion (Supine)}

\begin{itemize}
  \item Patient positioned supine; stabilize lower extremities and pelvis.
  \item Cue trunk flexion: patient raises upper trunk until inferior scapulae clear the table.
  \item Grade as follows:
    \begin{itemize}
      \item \textbf{5} – Hands behind head, full ROM.
      \item \textbf{4} – Arms crossed over chest.
      \item \textbf{3+} – Arms at sides, scapulae clear.
      \item \textbf{3} – Arms at sides, scapulae remain on table.
      \item \textbf{2} – Head and cervical spine only.
      \item \textbf{1–0} – Palpable contraction/no contraction.
    \end{itemize}
  \item Emphasize controlled movement and avoidance of substitution.
\end{itemize}
\end{labactivitybox}

\begin{labactivitybox}
\textbf{Activity: Core Control – Trunk Flexion Leg Lowering}

\begin{itemize}
  \item Patient supine with hips at 90°; arms crossed or hands by side.
  \item Patient slowly lowers legs while therapist monitors lumbar control.
  \item Grade based on angle at which control is lost (lumbar extension):
    \begin{itemize}
      \item \textbf{5} – Maintains control to table.
      \item \textbf{4} – Control lost past 30°.
      \item \textbf{3} – Control lost past 70°.
      \item \textbf{2} – Control lost above 80°.
      \item \textbf{1–0} – Palpable contraction/no control.
    \end{itemize}
  \item Document angle of control loss and correlate with trunk strength and motor control.
\end{itemize}
\end{labactivitybox}

\begin{labactivitybox}
\textbf{Activity: MMT – Trunk Extension (Prone)}

\begin{itemize}
  \item Patient prone; stabilize posterior thigh region.
  \item Cue patient to lift trunk until sternum clears the table.
  \item Grade as follows:
    \begin{itemize}
      \item \textbf{5} – Full ROM, hands behind head.
      \item \textbf{4} – Full ROM, arms folded behind back.
      \item \textbf{3} – Full ROM, arms at sides.
      \item \textbf{2} – Partial ROM.
      \item \textbf{1–0} – Palpable contraction/no movement.
    \end{itemize}
  \item Watch for compensatory motion or excessive lumbar lordosis.
\end{itemize}
\end{labactivitybox}

\section*{Rehabilitative Ultrasound Imaging – TrA and Multifidus}

Rehabilitative Ultrasound Imaging (RUSI) has emerged as a tool in the assessment and training of deep trunk stabilizers. According to Teyhen et al. (2007)\footnote{Teyhen DS, et al. (2007). The use of ultrasound imaging in rehabilitation of the lumbopelvic region: a clinical perspective. \textit{JOSPT}}, RUSI enables real-time visualization of motor control strategies, neuromuscular recruitment, and patient-specific adaptations. It has been shown to increase motor learning and feedback when retraining muscles like the transversus abdominis (TrA) and lumbar multifidus—key players in segmental stabilization.

\begin{labactivitybox}
\textbf{Activity: RUSI of Transversus Abdominis (TrA)}

\begin{itemize}
  \item Position patient supine with knees bent; probe placed in short axis just lateral to the umbilicus, oriented transversely.
  \item Visualize three abdominal wall layers: external oblique, internal oblique, and TrA (deepest).
  \item Instruct patient in abdominal drawing-in maneuver (ADIM) without pelvic movement.
  \item Record muscle thickness at rest and during contraction.
  \item Note activation pattern, symmetry, and substitution (e.g., rectus dominance or breath holding).
  \item \textit{Interpretation:} TrA contraction should produce a visible lateral slide or thickening without compensation. Symmetrical activation and absence of superficial dominance indicate effective segmental stabilization. Use this scan to guide retraining of deep core activation strategies.
\end{itemize}
\end{labactivitybox}

\begin{labactivitybox}
\textbf{Activity: RUSI of Lumbar Multifidus}

\begin{itemize}
  \item Position patient prone with pillow under abdomen for lumbar neutral.
  \item Place probe in short axis at L4–L5 (identify spinous process and lamina).
  \item Visualize multifidus as the hypoechoic, striated band lateral to the spinous process.
  \item Instruct patient to lift contralateral arm slightly to cue lumbar stabilization.
  \item Compare left vs. right thickness and contraction visibility.
  \item \textit{Interpretation:} Multifidus activation should result in a localized increase in muscle thickness without substitution from erector spinae. Asymmetry, delayed contraction, or absent response may indicate neuromuscular inhibition commonly seen in low back pain populations.
\end{itemize}
\end{labactivitybox}

\begin{tcolorbox}[colback=gray!5, colframe=black, title=Global Movers vs. Local Stabilizers]

Trunk muscle function can be broadly categorized into global movers and local stabilizers. \textbf{Global muscles} (e.g., rectus abdominis, external obliques, erector spinae) produce large-scale trunk motion and torque. \textbf{Local stabilizers} (e.g., transversus abdominis, lumbar multifidus) contribute to intersegmental control and postural stability with minimal external movement.

This distinction is functionally relevant in both assessment and rehabilitation. Weakness or poor recruitment of local stabilizers may lead to overuse of global muscles, altered motor control, or recurrent pain patterns.

\textbf{RUSI} helps visualize these patterns by showing symmetry, timing, and contraction quality. For example, TrA activation during an abdominal drawing-in maneuver or multifidus contraction during contralateral arm lift can confirm deep motor engagement.

\textit{Clinical relevance:} This approach is consistent with APTA Clinical Practice Guidelines recommending motor control training and segmental stabilization for managing low back pain, particularly in patients with movement coordination impairments or recurrent symptoms.\footnote{George SZ et al. (2021). Interventions for the Management of Acute and Chronic Low Back Pain: A Clinical Practice Guideline from the APTA. \textit{JOSPT}, 51(11):CPG1–CPG60. Delitto A et al. (2012). Low Back Pain Clinical Practice Guidelines. \textit{JOSPT}, 42(4):A1–A57.}

\end{tcolorbox}

\section*{Clinical Integration}

\begin{itemize}
  \item How do trunk MMT results relate to observed postural or movement control deficits in SFMA top tier tests?
 \item How did your observations during MMT, leg-lowering, and RUSI help you differentiate between global trunk movers (e.g., rectus abdominis) and local stabilizers (e.g., TrA, multifidus)?
  \item How did leg-lowering and RUSI feedback contribute to your understanding of lumbopelvic control?
  \item What are the limitations of RUSI and MMT in assessing core function? How might motor control tests complement strength-based assessments?
\end{itemize}

\section*{Next Steps}

\begin{itemize}
  \item Continue practicing trunk MMTs with clear grading criteria and appropriate stabilization.
  \item Review the anatomy and innervation of local trunk stabilizers, particularly the deep abdominal wall and lumbar multifidi.
  \item Practice your technique for obtaining clear TrA and multifidus RUSI images with reliable contraction visualization.
  \item Preview upcoming Week 9 content on cervical spine and quarter screens, especially the neurological connections from trunk to limb.
\end{itemize}

\section*{Checklist Summary}
\begin{itemize}
  \item [$\square$] Applied SFMA Top Tier movement logic to identify lumbopelvic control deficits.
  \item [$\square$] Performed trunk flexion and extension MMT using standardized criteria.
  \item [$\square$] Executed leg-lowering test to assess trunk motor control with angle of loss recorded.
  \item [$\square$] Identified primary trunk muscles with associated nerve and root innervation.
  \item [$\square$] Acquired and interpreted ultrasound images of TrA and multifidus.
  \item [$\square$] Analyzed muscle recruitment quality, symmetry, and substitution patterns.
\end{itemize}